\documentclass{VUMIFSlides}
\usepackage{times}
\usepackage[T1]{fontenc}
\usepackage{color}
\usepackage{verbatim}
\usepackage{graphicx}
\usepackage{fancyvrb}
\usepackage{bm}
\usepackage{amsfonts}
\usepackage{float}
\usepackage{hyperref}
\usepackage{textpos} % FOR LOGO
\usepackage[english,lithuanian]{babel}
\usepackage{animate}
\usepackage{caption} 
\usepackage{subfig}

\hypersetup{
    pdfauthor=Modestas Gražys,
    linktoc=all,     %set to all if you want both sections and subsections linked
}

\title[]{Grafeno gavybos proceso modeliavimas ir atpažinimas taikant kompiuterinės regos sprendimus}
\subtitle{}
\author[Modestas Gražys]{\texorpdfstring{Autorius: Modestas Gražys\\ Darbo vadovas: doc. prof. Linas Petkevičius}{Autorius: Modestas Gražys, Darbo vadovas: assoc. prof. Linas Petkevičius}} 
\institute[VU-MIF]{Vilniaus Universitetas \\ Matematikos ir Informatikos fakultetas \\ Informatikos institutas}
\date[2023-01-23]{2023-01-23}

\bibliography{bib}

% Let's get started
\begin{document}

\begin{frame}
  \maketitle
\end{frame}

% Content
\begin{frame}[c]{Motivation}    
    {\bf  Some list}
    \begin{itemize}
        \item  XXX
    \end{itemize}
    \vspace{1cm}
    {\bf Some other list}
    \begin{itemize}
        \item  XXX
    \end{itemize}
\end{frame}
% Introduction
\section{RevisedGoals}

\begin{frame}[c]{Tikslas}
    Ištirti ir atskirai pritaikyti skaitmeninius komponentus, reikalingus vėliau sudarinėjant grafeno sintezės simuliacijos įrankio modelį.
    \\
    \vspace{1cm}
    {\bf  Uždaviniai:}
    \begin{itemize}
        \item Pagal uždavinio ir įrankių charakteristikas apdoroti turimus duomenis.
        \item Apmokyti CLIP (Contrastive Language–Image Pre-training) modelį.
        \item Ištirti CLIP modelio pritaikomumą - gebėjimą klasifikuoti grafeno sintezės rezultatų duomenis.
        \item Apmokyti difuzijos modelį - generatorių.
        \item Ištirti generatoriaus gebėjimą besąlygiškai generuoti vaizdus atitinkančius grafeno sintezės reakcijos rezultatus.
        \item Rasti tinkamus hiperparametrus.
    \end{itemize}
\end{frame}
% Revised Dataset
\section{Data set}

\begin{frame}[c]{Duomenys}
    \begin{minipage}{0.45\textwidth}
    Grafeno sintezės oksidatorių klasės:
        \begin{itemize}
            \item NH4)2S2O8
            \item K2S2O8
            \item CrO3
            \item NH4)2S2O8\_P2O5
            \item K2S2O8\_P2O5
            \item CrO3\_P2O5
        \end{itemize}
    \end{minipage}
    \begin{minipage}{0.5\textwidth}
        \begin{figure}
            \centering
            \includegraphics[width=1\textwidth]{img/examples.png}
            \caption{Grafito bisulfato (GBS) mėginių optinės mikroskopijos nuotraukos}
            \label{fig:graphene_data}
        \end{figure}
    \end{minipage}
\end{frame}

% Methods used
\section{Methods used}

\begin{frame}[c]{Some model}
    \begin{minipage}{0.45\textwidth}
        \begin{itemize}
            \item a..
        \end{itemize}
    \end{minipage}
    \begin{minipage}{0.5\textwidth}
        \centering
        % \includegraphics[width=1\textwidth]{imgs/XX.png}
    \end{minipage}
    Citation from bib.bib \cite{lecun2015deep}
\end{frame}

% Analysis results
\begin{frame}[c]{Įrankiai ir mokymas}
    \begin{minipage}[t]{0.65\textwidth}
        CLIP mokymas:
        \begin{itemize}
            \item Kryžminės entropijos nuostolių funkcija
            \item 200 epochų
            \item Adam optimizatorius
            \item Mokymo žingsnio kosinusinio
kitimo planuotojas (\textit{cosine learning rate scheduler})
        \end{itemize}
        Difuzerio mokymas:
        \begin{itemize}
            \item Vidutinės
kvadratinės paklaidos nuostolių funkcija
            \item 200 epochų
            \item AdamW optimizatorius
            \item Mokymo žingsnio kosinusinio
kitimo planuotojas (\textit{cosine learning rate scheduler})
        \end{itemize}
    \end{minipage}
    \begin{minipage}[t]{0.29\textwidth}
        \begin{figure}
            \centering
            \includegraphics[width=0.75\textwidth]{img/Google_Colaboratory_SVG_Logo.svg.png}
        \end{figure}
        \begin{figure}
            \centering
            \includegraphics[width=0.75\textwidth]{img/Pytorch_logo.png}
        \end{figure}
        \begin{figure}
            \centering
            \includegraphics[width=0.75\textwidth]{img/6308b84661b3e2a522f01468.png}
        \end{figure}
    \end{minipage}
\end{frame}
\begin{frame}[c]{Rezultatai}
    \begin{minipage}[t]{0.49\textwidth}
        \begin{itemize}
            \item CLIP modelis grafeno sintezės reakcijos rezultato nuotraukas klasifikuoja apytiksliai 16.67\% įsitikinimo (\textit{confidence}) lygmeniu.
            \item Difuzijos modelio optimali mokymo žingsnio reikšmė $3 * 10^{-4}$.
            \item Difuzijos modelio optimali mokymo žingsnio planuotojo     apšilimo reikšmė 0\%.
            \item Difuzijos modelio generuojamos nuotraukos panašios į mokymo duomenų nuotraukas.
        \end{itemize}
    \end{minipage}
    \begin{minipage}[t]{0.49\textwidth}
       \begin{figure}
            \includegraphics[width=1\textwidth]{img/graphene-diffuser-inference.png}
            \caption{Generatyviniu difuzijos modeliu sugeneruoti grafeno vaizdai}
        \end{figure}
    \end{minipage}
\end{frame}

\begin{comment}
\begin{frame}[c]{Visualization I}
    \begin{figure}
        \includegraphics[width=1\textwidth]{img/sample1.png}
    \end{figure}
\end{frame}

\begin{frame}[c]{Visualization II}
    \begin{figure}
        \begin{tikzpicture}
            \node[inner sep=0pt] (A) at (0,0){\includegraphics[width=.25\textwidth]{img/sample1.png}};
            \node[inner sep=0pt] (B) at (5,-4){\includegraphics[width=.25\textwidth]{img/sample1.png}};
            \draw[<->,thick] (A.south east) -- (B.north west)node[midway,fill=white] {Some arrow};
            \node (A) at (6,0) {\(\displaystyle\left|
                \frac{
                    \left(\sum_{i=1}^n a_i^2\right) \cdot
                    \left(\sum_{i=1}^n b_i^2\right)
                }{ (\sum_{i=1}^n \underbrace{a_ib_i}_{\subnode{brace}{}})^2} \right| \geq 1\)};
        \end{tikzpicture}   
    \end{figure}        
\end{frame}
   
\begin{frame}[c]{Visualization III}    
    \begin{minipage}[t]{0.49\textwidth}
        \begin{itemize}
            \item XXX
            \item XXX
            \item XXX
        \end{itemize}
    \end{minipage}
    \begin{minipage}[t]{0.49\textwidth}
       \begin{figure}
            \includegraphics[width=1\textwidth]{img/sample1.png}
        \end{figure}
    \end{minipage}
\end{frame}


\begin{frame}[fragile]{Visualization IV}
    \begin{figure}[h!] 
        \begin{animateinline}[autoplay,loop]{6}
            \multiframe{3}{i=1+1}{	% {upper_index}{iterate + 1}
                \begin{tikzpicture}[scale = 0.99]
	 	     %\pgfmathsetmacro\c{\i+0} % c can be new variable of indexing
                    \node at (-8, 0) {\i};
		          \node[inner sep=0pt] (A) at (0,0)
                    {\includegraphics[width=.75\textwidth]{img/sample\i.png}};
                \end{tikzpicture}%   
            }
        \end{animateinline}
    \end{figure}
\end{frame}
\end{comment}
% Benefits / conclusion
\section{Benefits}

 \begin{frame}[c]{Išvados}
    \begin{itemize}
        \item Pagal dabartinį projektavimą CLIP modelis nesugeba išmokti užtikrintai klasifikuoti grafeno sintezės vaizdų pagal oksidatoriaus klases. Visiems klasės spėjimo variantams modelis priskiria apylygias tikimybes.
        \item Difuzijos modelis (generatorius) generuoja panašius vaizdus į tuos, kokių būtų tikimąsi, tačiau šie vaizdai nėra pakankamai detalūs - detaliai aiškūs.
        \item Remiantis paveikslėlių generavimo iš tekstinės įvesties užduotimis (\textit{text2image}), šiuos abu modelius galima adaptuoti grafeno sintezės reakcijų rezultatų prognozavimui.
    \end{itemize}
\end{frame}

% Bibliography
\begin{frame}[allowframebreaks]{Literatūra}
    \printbibliography 
\end{frame}
 
% \appendix

\end{document}
